\documentclass[a4paper]{article}
\usepackage[english]{babel}
\usepackage{fancyhdr}
\usepackage{graphicx}
\usepackage[latin1]{inputenc}
\usepackage{lastpage}
\usepackage{longtable}
\usepackage{multirow}
\usepackage{rotating}
\usepackage{tabls} \setlength\tablinesep{3pt}
\usepackage{url}

\setlength{\parindent}{0mm}
\renewcommand{\familydefault}{\sfdefault}

\begin{document}
\begin{titlepage}
	\bigskip
	\bigskip
	\bigskip
	\begin{center}
		\includegraphics{images/informatik-tuwien.png}\\
		\vspace{1cm}
		\Huge{\textsf{\textbf{Design Methods for Distributed Systems\\}}}
		\vspace{1cm}
		\huge{\textsf{\textbf{Architecture\\}}}
		\vspace{1cm}
		\LARGE{\textsf{VU Entwurfsmethoden f\"{u}r Verteilte Systeme\\}}
		\vspace{1cm}
		\Large{\textsf{SS 2008\\}}
	\end{center}
	\bigskip
	\begin{center}
		\large{\textsf{\textbf{Group 69}}}
	\end{center}
	\begin{tabular}{llll}
	\textsf{Martina Lindorfer} & \textsf{0626770} & \textsf{066937} & \url{e0626770@student.tuwien.ac.at}\\
	\textsf{Florian Lukavsky} & \textsf{0325558} & \textsf{066937} & \url{e0325558@student.tuwien.ac.at}\\
	\textsf{Daniel Priewasser} & \textsf{0626777} & \textsf{066937} & \url{e0626777@student.tuwien.ac.at}\\
	\textsf{Gerald Scharitzer} & \textsf{0127228} & \textsf{066937} & \url{e0127228@student.tuwien.ac.at}\\
	\textsf{Dirk Wallerstorfer} & \textsf{0626775} & \textsf{066937} & \url{e0626775@student.tuwien.ac.at}\\
	\end{tabular}
\end{titlepage}

\clearpage

\pagestyle{fancy}				
\pagenumbering{arabic}				
\cfoot{\small{\thepage\ / \pageref{LastPage}}}

\tableofcontents

\newpage
\section{Logical View}

\cite{zdun2008}

\subsection{Common}
\subsubsection{Marshaller}

The marshaller translates objects into sequences of bytes and vice versa.

\subsubsection{Interface Description}

The interface description is the source for generating the interface for client proxies and remote objects.

\subsection{Server}

\subsubsection{Remote Object}

\subsubsection{Invoker}

The invoker calls the marshaller to translate request messages to invocations.

\subsubsection{Server Request Handler}

The server request handler receives request messages and dispatches them to the invokers.
The server request handler sends response messages to the client request handlers.

\subsubsection{Server Protocol Plugin}

\subsection{Client}

\subsubsection{Client Protocol Plugin}

\subsubsection{Client Request Handler}

\subsubsection{Requestor}

\subsubsection{Client Proxy}

\subsection{Protocols}

\subsubsection{Communication Protocols}

\paragraph{Minimal Communication Protocol}

The minimal communication protocol has no protocol header and therefore transmits byte streams without interpreting them.
The marshalling protocol is determined by the communication end points, which process only single protocols.
The sender and receiver are determined by the communication end points.

\paragraph{Typed Communication Protocol}

The typed communication protocol prefixes the message body with a 4 byte header, which specifies the protocol of the message body.
The sender and receiver are determined by the communication end points.

\subsubsection{Marshalling Protocols}

\paragraph{Binary Serialization Protocol}

The invocation is analysed and serialized via reflection.

\paragraph{Java Serialization Protocol}

The invocation is serialized by the default Java object serialization.

\paragraph{Axi2 Java Web Service Protocol}

The marshaller translates invocations to SOAP messages and vice versa.

\subsection{Configuration}

\subsubsection{Properties}

peer.local.host.name\\
peer.local.ip.address\\
peer.local.port\\

\section{Development View}

\subsection{Tools}

	\begin{tabular}{|l|l|l|}
	\hline
	Development Environment & Eclipse 3.3.2 & \url{http://www.eclipse.org/}\\
	\hline
	Programming Language & Java 5 & \url{http://java.sun.com/}\\
	\hline
	Web Services & Axis2 & \url{http://ws.apache.org/axis2/}\\
	\hline
	Unit Tests & JUnit 4 & \url{http://junit.org/}\\
	\hline
	Build Tool & Ant 1.7 & \url{http://ant.apache.org/}\\
	\hline
	Version Control & Subclipse 1.2.4 & \url{http://subclipse.tigris.org/}\\
	\hline
	\end{tabular}

\subsection{Naming Rules}

The top level package is "evs" and there is only one level of sub packages.\\
Interface names start with "I".\\
Abstract class names start with "A".\\

\subsection{Packages}

	\begin{tabular}{|l|l|}
	\hline
	evs.comm & communcation framework\\
	\hline
	evs.entities & business entities\\
	\hline
	evs.junit & JUnit tests\\
	\hline
	\end{tabular}

\section{Process View}

Communication is based on Java sockets.

\section{Physical View}

\section{Scenarios}

\subsection{Use Cases}

\subsubsection{Create Customer}

\subsubsection{Create Product}

\subsubsection{Find Product}

\subsubsection{Buy Product}

\section{Rationale}

\section{Extension 1}

\subsection{Asynchronous Invocation}
\subsubsection{Fire and Forget}
\subsubsection{Poll Object}
\subsubsection{Result Callback}

\subsection{Interceptor Architecture}
\subsubsection{Extension Patterns}

\subsection{Activation of Distributed Objects}
\subsubsection{Life Cycle Management Patterns}

\subsection{Federation}
\subsubsection{Configuration Groups}

\section{Extension 2}
\subsection{Dynamic Deployment}
\subsubsection{Command}
\subsubsection{Command Language}
\subsubsection{Object System Layer}

\newpage
\bibliographystyle{plain}
\bibliography{bibliography}

\end{document}
